\documentclass[10pt,a3paper]{article}
\usepackage[margin=10mm]{geometry}
\usepackage[T1]{fontenc}
\usepackage[utf8]{inputenc}
\usepackage[french]{babel}
\usepackage{lmodern}
\usepackage{microtype}
\makeatletter\def\input@path{{../../../Style/}}\makeatother
\NeedsTeXFormat{LaTeX2e}
\ProvidesPackage{TLPosterCarteMentale}[2026/08/03 Posters et cartes mentales SI]

\RequirePackage{tikz}
\RequirePackage[most]{tcolorbox}
\RequirePackage{xcolor}
\RequirePackage{amsmath,amssymb}
\RequirePackage{graphicx}
\usetikzlibrary{positioning,arrows.meta,calc,shapes.geometric,mindmap,backgrounds,fit}

\definecolor{TLPosterBlue}{RGB}{24,86,141}
\definecolor{TLPosterLight}{RGB}{234,243,250}
\definecolor{TLPosterAccent}{RGB}{232,126,34}
\definecolor{TLPosterGreen}{RGB}{52,152,102}
\definecolor{TLPosterRed}{RGB}{192,57,43}
\definecolor{TLPosterGray}{RGB}{80,90,100}

\newcommand{\TLPosterTitre}[2]{%
  \begin{tcolorbox}[enhanced,boxrule=0pt,arc=0pt,colback=TLPosterBlue,left=6mm,right=6mm,top=4mm,bottom=4mm]
    {\color{white}\bfseries\LARGE #1}\hfill{\color{white}\large #2}
  \end{tcolorbox}%
}

\newtcolorbox{TLPosterBloc}[2][]{
  enhanced,breakable,
  colback=TLPosterLight,
  colframe=TLPosterBlue,
  colbacktitle=TLPosterBlue,
  coltitle=white,
  fonttitle=\bfseries,
  title={#2},
  boxrule=0.8pt,
  arc=2mm,
  left=3mm,right=3mm,top=2mm,bottom=2mm,
  #1
}

\newtcolorbox{TLPosterFormule}[1][]{
  enhanced,
  colback=white,
  colframe=TLPosterAccent,
  boxrule=1pt,
  arc=2mm,
  fontupper=\bfseries,
  halign=center,
  #1
}

\newcommand{\TLPosterMethode}[1]{%
  \begin{tcolorbox}[enhanced,colback=TLPosterGreen!8,colframe=TLPosterGreen,title=Méthode,fonttitle=\bfseries]
  #1
  \end{tcolorbox}%
}

\newcommand{\TLPosterAttention}[1]{%
  \begin{tcolorbox}[enhanced,colback=TLPosterRed!6,colframe=TLPosterRed,title=Point de vigilance,fonttitle=\bfseries]
  #1
  \end{tcolorbox}%
}

\tikzset{
  TLmindroot/.style={concept,concept color=TLPosterBlue,text=white,font=\bfseries\large,minimum size=34mm},
  TLmindlevel1/.style={level 1 concept/.append style={font=\bfseries,minimum size=23mm,sibling angle=45}},
  TLmindlevel2/.style={level 2 concept/.append style={font=\small,minimum size=17mm}},
  TLarrow/.style={-{Stealth[length=3mm]},thick,draw=TLPosterBlue},
  TLbox/.style={rounded corners=2mm,draw=TLPosterBlue,fill=TLPosterLight,align=center,inner sep=3mm}
}

\newenvironment{TLCarteMentale}[1]{%
  \begin{tikzpicture}[mindmap,grow cyclic,concept color=TLPosterBlue,every node/.style={concept},TLmindlevel1,TLmindlevel2]
  \node[TLmindroot]{#1}
}{;
  \end{tikzpicture}
}

\newcommand{\TLBranche}[3][]{child[concept color=#2]{node[#1]{#3}}}

\endinput

\pagestyle{empty}
\renewcommand{\familydefault}{\sfdefault}
\begin{document}
\TLPosterCourseHeader{1}{Ingénierie système}{Besoin, exigences, SysML et architecture fonctionnelle}
\vspace{2mm}
\begin{minipage}[t]{0.61\linewidth}
\begin{TLPosterSavoirs}
\begin{itemize}
  \item définir système, frontière, environnement, parties prenantes et phases du cycle de vie
  \item distinguer système souhaité, simulé et réel et interpréter les écarts de performances
  \item exprimer le besoin par les services attendus et le contexte d'utilisation
  \item formuler une exigence \textbf{MUST} : mesurable, utile, simple et traçable/testable
  \item choisir le diagramme SysML adapté : Use Case, req, BDD, IBD, séquence, activité ou états
\end{itemize}
\end{TLPosterSavoirs}
\end{minipage}\hfill
\begin{minipage}[t]{0.36\linewidth}
\TLPosterKey{Quel diagramme ?}{Besoin : Use Case / contexte\\ Exigences : req\\ Composition : BDD\\ Interfaces et flux : IBD\\ Scénario : séquence\\ Modes : états}
\end{minipage}
\vfill
\begin{minipage}[t]{0.49\linewidth}
\begin{TLPosterBloc}[Si on demande le besoin ou le contexte]{TLPosterBlue}
Identifier la \textbf{phase de vie}, les \textbf{parties prenantes} et les éléments extérieurs en interaction. Le Use Case exprime les services rendus ; le BDD de contexte fixe la frontière et les interacteurs. Ne pas proposer de solution technique à ce stade.
\end{TLPosterBloc}
\begin{TLPosterBloc}[Si on demande les exigences]{TLPosterPurple}
Transformer le besoin en exigences vérifiables. Pour chaque exigence préciser au minimum une grandeur/critère, un niveau attendu et, si possible, le moyen de vérification. Le diagramme \textbf{req} assure la hiérarchie et la traçabilité jusqu'aux éléments qui satisfont ou vérifient l'exigence.
\end{TLPosterBloc}
\end{minipage}\hfill
\begin{minipage}[t]{0.49\linewidth}
\begin{TLPosterBloc}[Si on demande l'architecture]{TLPosterGreen}
Utiliser le \textbf{BDD} pour répondre à « de quoi le système est-il constitué ? ». Utiliser l'\textbf{IBD} pour répondre à « qui échange quoi avec qui ? » : matière, énergie ou information. Relier ensuite fonctions, chaîne d'information, chaîne de puissance et solutions.
\end{TLPosterBloc}
\begin{TLPosterBloc}[Si on demande le fonctionnement]{TLPosterTeal}
Séquence : échanges ordonnés dans le temps. Activité : enchaînement d'actions et flux. États : modes, événements, gardes et transitions. Toujours choisir le diagramme selon la question posée plutôt que de représenter tout le système.
\end{TLPosterBloc}
\end{minipage}
\vfill
\begin{TLPosterMethode}
\begin{enumerate}
  \item Lire la question et identifier le \textbf{point de vue demandé} : besoin, exigence, structure ou comportement.
  \item Fixer la frontière du système et la phase du cycle de vie étudiée.
  \item Relever parties prenantes, interacteurs, flux et services sans inventer de solution.
  \item Choisir le diagramme SysML qui répond directement à la question.
  \item Vérifier la cohérence : chaque exigence doit être mesurable et reliée à un élément du modèle.
  \item Comparer souhaité, simulé et réel pour conclure sur vérification puis validation du besoin.
\end{enumerate}
\end{TLPosterMethode}
\vfill
\TLPosterRetenir{partir du besoin, choisir le bon point de vue SysML, tracer les exigences jusqu'aux solutions puis vérifier et valider.}
\end{document}
